% =============================================================================
% 060 / 068
% Status: raw
% =============================================================================
% Notes:
%
% Source question:
% 新冠肺炎疫情的衝擊，牆內海外無一倖免，
% 這對吳語復興有很嚴重的影響了吧？
% 時間線更加緊迫了。
% 
% 祝早日順利建構好、階段性完成部分辭典，
% 祝現有的語言記憶和文物儲存良好。
% =============================================================================

\chapter[新冠肺炎疫情的衝擊，牆內海外無一倖免， 這對吳語復興有很嚴重的影響了吧？ 時間線更加緊迫了。 祝早日順利建構好、階段性完成部分辭典， 祝現有的語言記憶和文物儲存良好。]{新冠肺炎疫情的衝擊，牆內海外無一倖免， \\[0.35em] 這對吳語復興有很嚴重的影響了吧？ \\[0.35em] 時間線更加緊迫了。 \\[0.35em] 祝早日順利建構好、階段性完成部分辭典， \\[0.35em] 祝現有的語言記憶和文物儲存良好。}
\qaanswer{
吳語播音體被世人嚴肅的對待，並被推廣學習極大可能是在吳語區現有吳語(包括吳音藍青話)消亡以後的事了。只有擁有文字紀錄的語言才有可能做為標本留存下去，而在未來，吳語播音體可能是唯一可信的證明吳語曾經存在過，並區別於Mandarin的證據了。而未來在吳語業已消亡的很長時間之後，如果還能復活，很有可能靠的是現在構建中的基於梳理原生吳語後的吳語播音體的體系。我們所要做的就是為未來保留這種可能性。
}

